% Fiche VHDL — recto-verso A4 paysage
% Source : VHDLCheatSheet.md (I. El Badawy) — repris, traduit, corrigé et complété
%   (std_logic, numeric_std, conversions, paquets, fonctions, FSM, compteur, RAM,
%   testbench complet, pièges).
% Sœur de ../Verilog/verilog_cheat.tex (même gabarit).
% Compiler : make (ici ou à la racine), ou xelatex vhdl_cheat.tex — une seule passe suffit
% Apparence : ../librairie (fiche.sty commun + thème). Macros de contenu :
%   colonne (3 par page, séparées par \colbreak) · bloc{titre}{cat}
%   catégories : syn typ ctl ex io lib (couleurs dans le thème)
%   code (listing numéroté) · \V|...| code en ligne · «...» dans le code = à remplir
%   \sub sous-titre · \note remarque · \piege piège · \badge{TXT}
%   optab[l1] étiquette | texte · tab3[l1]{l2} étiquette | code | remarque
% Le code est en police à chasse fixe : les commentaires sont alignés à la main
% (attention : les « » ne sont pas affichés, ils ne comptent pas dans l'alignement).
% Après compilation : grep "\[fiche\]" vhdl_cheat.log donne l'espace libre par colonne.
\documentclass{article}
\usepackage[vhdl]{../librairie/fiche}   % apparence : ../librairie/fiche.sty + theme-vhdl.tex
% --- Langage VHDL pour listings (insensible à la casse) ---
% Classes de mots-clés : [1] rose · [2] cyan · [3] vert · [4] couleur du thème
% (styles dans ../librairie/fiche.sty). Chargé par vhdl_cheat.tex.
% [1] mots-clés · [2] types · [3] fonctions et paquets
% Pas de chaînes '...' : l'apostrophe sert aussi aux attributs (x'length).
\lstdefinelanguage{CODE}{
  sensitive=false,
  morekeywords=[1]{library,use,all,entity,is,port,generic,map,end,architecture,of,
    begin,signal,variable,constant,component,process,if,then,elsif,else,case,when,
    others,for,loop,while,generate,in,out,inout,buffer,downto,to,wait,until,on,
    report,severity,assert,after,function,procedure,return,type,subtype,array,record,
    range,package,body,select,with,null,exit,next,and,or,not,nand,nor,xor,xnor,
    sll,srl,sla,sra,rol,ror,mod,rem,abs,open,alias,attribute,work,impure,pure},
  morekeywords=[2]{std_logic,std_logic_vector,std_ulogic,std_ulogic_vector,unsigned,
    signed,integer,natural,positive,boolean,bit,bit_vector,real,time,string,character},
  morekeywords=[3]{ieee,std_logic_1164,numeric_std,rising_edge,falling_edge,
    to_unsigned,to_signed,to_integer,resize,shift_left,shift_right,rotate_left,
    rotate_right,now,note,warning,error,failure,std,env,stop,finish,to_hstring,image},
  morecomment=[l]{--}, morestring=[b]",
}
   % coloration syntaxique : \lstdefinelanguage{CODE}

\begin{document}
%%%%%%%%%%%%%%%%%%%%%%%%%%%%%%%%%%%%%%%%%%%%%%%%%%%%%%%%%%%%%%%%%%%%%%%%%%%%%%
%                         RECTO : LE LANGAGE
%%%%%%%%%%%%%%%%%%%%%%%%%%%%%%%%%%%%%%%%%%%%%%%%%%%%%%%%%%%%%%%%%%%%%%%%%%%%%%
\entete{VHDL}{recto · le langage}
\noindent
% ======================= colonne 1 : structure =======================
\begin{colonne}
\begin{bloc}{Bibliothèques}{lib}
\begin{code}
library ieee;
use ieee.std_logic_1164.all;  -- std_logic, rising_edge
use ieee.numeric_std.all;     -- unsigned, signed, +, -
use work.outils.all;          -- un paquet du projet
\end{code}
\piege{jamais \V|std_logic_arith| ni \V|std_logic_unsigned| : non standard, en conflit avec \V|numeric_std|.}
\end{bloc}

\begin{bloc}{Entité : l'interface}{syn}
\begin{code}
entity compteur is
  generic (
    N    : positive := 8;     -- paramètre
    MAXV : natural  := 255
  );
  port (
    clk, rst : in  std_logic;
    en       : in  std_logic;
    q        : out std_logic_vector(N-1 downto 0);
    pleine   : out std_logic  -- pas de ; final
  );
end entity;
\end{code}
\begin{optab}[12mm]
\V|in|     & lecture seule \\
\V|out|    & écriture (lisible en interne en VHDL-2008) \\
\V|inout|  & bidirectionnel (bus trois états) \\
\V|buffer| & sortie relisible (à éviter)
\end{optab}
\note{Ports et génériques sont séparés par des \V|;|, pas des virgules.}
\end{bloc}

\begin{bloc}{Architecture : le contenu}{syn}
\begin{code}
architecture rtl of compteur is
  -- zone déclarative : signaux, constantes, types
  signal cnt    : unsigned(N-1 downto 0);
  signal data_b : std_logic_vector(7 downto 0);
  constant ZERO : unsigned(cnt'range) := (others => '0');
begin
  -- instructions concurrentes (ordre indifférent)
  q      <= std_logic_vector(cnt);
  pleine <= '1' when cnt = MAXV else '0';
end architecture;
\end{code}
\note{Plusieurs architectures par entité possibles ; \V|rtl|, \V|behav| sont des noms libres.}
\end{bloc}

\begin{bloc}{Paquet et fonction}{lib}
\begin{code}
package outils is       -- déclarations
  constant LARG : positive := 32;
  function maxi(a, b : integer) return integer;
end package;
package body outils is  -- définitions
  function maxi(a, b : integer) return integer is
  begin
    if a > b then return a; else return b; end if;
  end function;
end package body;
\end{code}
\note{Une fonction rend une valeur, sans \V|wait|. Une \V|procedure| peut attendre (\V|wait|) et a des paramètres \V|out|.}
\end{bloc}
\end{colonne}\colbreak
% ======================= colonne 2 : types =======================
\begin{colonne}
\begin{bloc}{std\_logic : 9 valeurs}{typ}
\begin{tab3}[11.5mm]{21.5mm}
\V|'U'|         & non initialisé  & début de simulation \\
\V|'X'|         & conflit fort    & deux sources opposées \\
\V|'0'| \V|'1'| & niveaux forts   & valeurs logiques \\
\V|'Z'|         & flottant        & haute impédance (trois états) \\
\V|'W'|         & conflit faible  & \\
\V|'L'| \V|'H'| & 0 et 1 faibles  & pull-down, pull-up \\
\V|'-'|         & indifférent     & optimisation de synthèse
\end{tab3}
\note{Type résolu : plusieurs pilotes se combinent (\V|'0'| et \V|'1'| donnent \V|'X'|).}
\end{bloc}

\begin{bloc}{Types}{typ}
\begin{optab}[31.5mm]
\V|std_logic_vector|     & bus de bits, sans valeur numérique \\
\V|unsigned|, \V|signed| & bus avec arithmétique (\V|numeric_std|) \\
\V|integer range 0 to 9| & 32 bits si pas de \V|range| \\
\V|natural|, \V|positive| & entiers $\geq 0$, $\geq 1$ \\
\V|boolean|              & résultat des comparaisons \\
\V|time|                 & \V|ns|, \V|us|... (simulation)
\end{optab}
\begin{code}
type etat_t is (IDLE, LIRE, ECRIRE);  -- énuméré (FSM)
subtype octet is std_logic_vector(7 downto 0);
type mem_t is array (0 to 255)
  of std_logic_vector(7 downto 0);    -- tableau
type pix_t is record
  r, g, b : unsigned(7 downto 0);
end record;
signal mem : mem_t := (others => (others => '0'));
\end{code}
\end{bloc}

\begin{bloc}{Conversions (numeric\_std)}{lib}
\par\vspace{1mm}\centering
\begin{tikzpicture}[x=1mm,y=1mm]
  % chaîne verticale : descendre à gauche, remonter à droite
  \node[ty, minimum width=30mm] (v) at (0,0)   {std\_logic\_vector};
  \node[ty, minimum width=30mm] (u) at (0,-15) {unsigned \textcolor{ink!50}{|} signed};
  \node[ty, minimum width=30mm] (i) at (0,-30) {integer};
  \draw[cv] (-6,-2.6)  -- node[fn,left,align=right]{unsigned(v)\\signed(v)}               (-6,-12.4);
  \draw[cv] (6,-12.4)  -- node[fn,right,align=left]{std\_logic\_vector(u)\\std\_logic\_vector(s)} (6,-2.6);
  \draw[cv] (-6,-17.6) -- node[fn,left,align=right]{to\_integer(u)\\to\_integer(s)}       (-6,-27.4);
  \draw[cv] (6,-27.4)  -- node[fn,right,align=left]{to\_unsigned(i, N)\\to\_signed(i, N)} (6,-17.6);
  \node[fn,anchor=east,align=right] at (u.west) {signed(u), unsigned(s)\\(mêmes bits)~};
  \node[fn,anchor=west,align=left,text=ink!70] at (u.east) {~resize(x, N)\\(change la taille)};
\end{tikzpicture}\par\raggedright
\begin{code}
y <= std_logic_vector(unsigned(a) + unsigned(b));
n <= to_integer(unsigned(sel));  -- index
r <= resize(s, 16);              -- signe étendu
\end{code}
\end{bloc}

\begin{bloc}{Attributs}{typ}
\begin{optab}[21mm]
\V|v'length|         & nombre d'éléments \\
\V|v'range|           & \V|7 downto 0| : pour \V|for i in v'range| \\
\V|v'high|, \V|v'low| & bornes \\
\V|v'left|, \V|v'right| & bornes dans l'ordre déclaré \\
\V|T'image(x)| & texte de \V|x| (\V|integer'image(n)|) \\
\V|clk'event|         & front : préférer \V|rising_edge(clk)|
\end{optab}
\end{bloc}
\end{colonne}\colbreak
% ======================= colonne 3 : instructions =======================
\begin{colonne}
\begin{bloc}{Opérateurs}{syn}
\begin{tab3}[13mm]{28mm}
Arithm.  & \V|+ - * / **| & sur \V|unsigned|, \V|signed|, \V|integer| \\
         & \V|mod rem abs| & \V|mod| du signe du diviseur, \V|rem| du dividende \\
Logiques & \V|and or not xor| & aussi \V|nand nor xnor| ; bit à bit sur vecteurs \\
Relation & \V|= /= < <= > >=| & \V|/=| différent ; rendent un \V|boolean| \\
Décalage & \V|shift_left(u, 2)| & \V|shift_right| étend le signe sur \V|signed| \\
Concat.  & \V|a & b| & assemble des vecteurs \\
Affect.  & \V|<=| / \V|:=| & signal / variable
\end{tab3}
\piege{\V|and| et \V|or| ont la même priorité : \V|a and b or c| est refusé, parenthéser.}
\end{bloc}

\begin{bloc}{Instructions concurrentes}{ctl}
\begin{code}
y <= "0001" when a = "00" else  -- conditionnelle
     "0010" when a = "01" else
     "0100" when a = "10" else
     "1000";                    -- else final
with sel select                 -- sélective
  m <= x0 when "00",
       x1 when "01",
       x2 when others;
g : for i in 0 to N-1 generate  -- répliquer
  y(i) <= a(i) xor b(i);
end generate;
\end{code}
\end{bloc}

\begin{bloc}{Process : le séquentiel}{ctl}
\begin{code}
process (clk, rst)                -- liste de sensibilité
  variable tmp : integer := 0;    -- locale au process
begin
  if rst = '1' then
    ...
  elsif rising_edge(clk) then
    case op is
      when "00"         => tmp := tmp + 1;
      when "01" | "10"  => tmp := 0;
      when others       => null;  -- others obligatoire
    end case;
    for i in v'range loop r(i) <= v(7 - i); end loop;
  end if;
end process;
\end{code}
\note{Signal : nouvelle valeur à la fin du process ; variable \V|:=| : tout de suite.}
\end{bloc}

\begin{bloc}{Littéraux et agrégats}{syn}
\begin{code}
'0' '1' 'Z'                -- std_logic : apostrophes
"1010"   x"FF"   o"17"     -- vecteurs : guillemets
16#FF#   2#1010#   1_000_000  -- entiers dans une base
10 ns   1.5 us             -- time : espace obligatoire
(others => '0')            -- tout à 0, toute taille
(0 => '1', others => '0')  -- bit 0 à 1
a & b      v(7 downto 4)   -- concaténation, tranche
\end{code}
\end{bloc}
\end{colonne}

%%%%%%%%%%%%%%%%%%%%%%%%%%%%%%%%%%%%%%%%%%%%%%%%%%%%%%%%%%%%%%%%%%%%%%%%%%%%%%
%                VERSO : CONCEPTION ET SIMULATION
%%%%%%%%%%%%%%%%%%%%%%%%%%%%%%%%%%%%%%%%%%%%%%%%%%%%%%%%%%%%%%%%%%%%%%%%%%%%%%
\newpage
\entete{VHDL}{verso · conception et simulation}
\noindent
% ======================= colonne 1 : circuits =======================
\begin{colonne}
\begin{bloc}{Logique combinatoire}{ex}
\begin{code}
process (all)  -- VHDL-2008 : tout lu
begin
  y <= '0';    -- défaut : pas de latch
  case sel is
    when "00"   => y <= a;
    when "01"   => y <= b;
    when others => null;
  end case;
end process;
-- additionneur complet (flot de données)
sum  <= a xor b xor cin;
cout <= (a and b) or (cin and (a xor b));
\end{code}
\note{Avant 2008 : mettre dans la liste de sensibilité tous les signaux lus, sinon simulation et synthèse divergent.}
\end{bloc}

\begin{bloc}{Registre N bits}{ex}
\begin{code}
-- generic (N : positive := 16) ; d, q : N bits
process (clk, rst)    -- reset asynchrone
begin
  if rst = '1' then
    q <= (others => '0');
  elsif rising_edge(clk) then
    if en = '1' then  -- clock enable
      q <= d;
    end if;
  end if;
end process;
\end{code}
\note{Reset synchrone : \V|process (clk)| et tester \V|rst| à l'intérieur de \V|rising_edge(clk)|.}
\end{bloc}

\begin{bloc}{Compteur}{ex}
\begin{code}
process (clk)
begin
  if rising_edge(clk) then
    if rst = '1' then       -- reset synchrone
      cnt <= (others => '0');
    elsif en = '1' then
      if cnt = MAXV then cnt <= (others => '0');
      else cnt <= cnt + 1;  -- unsigned + integer
      end if;
    end if;
  end if;
end process;
\end{code}
\end{bloc}

\begin{bloc}{Décalage et détection de front}{ex}
\begin{code}
process (clk)
begin
  if rising_edge(clk) then
    sr   <= sr(6 downto 0) & din;  -- registre à décalage
    b_r  <= b;           -- b retardé d'un cycle
  end if;
end process;
front <= b and not b_r;  -- impulsion 0 -> 1
\end{code}
\end{bloc}

\end{colonne}\colbreak
% ======================= colonne 2 : hiérarchie, FSM =======================
\begin{colonne}
\begin{bloc}{Instancier un composant}{lib}
\begin{code}
-- direct (VHDL-93) : pas de déclaration à faire
u_reg : entity work.reg
  generic map (N => 8)
  port map (clk => clk, rst => rst, en => '1',
            d => d_in, q => open);  -- open : non relié
-- classique : déclarer le composant, puis
component fa is
  port (a, b, cin : in  std_logic;
        sum, cout : out std_logic);
end component;
-- association par position (ordre des ports)
f0 : fa port map (a(0), b(0), cin, s(0), c(0));
\end{code}
\note{Préférer l'association par nom (\V|=>|).}
\end{bloc}

\begin{bloc}{Additionneur N bits (generate)}{lib}
\begin{code}
signal c : std_logic_vector(N downto 0);
...
c(0) <= cin;
g : for i in 0 to N-1 generate
  fa_i : entity work.fa
    port map (a => a(i), b => b(i), cin => c(i),
              sum => s(i), cout => c(i+1));
end generate;
cout <= c(N);
-- plus simple, sans composant :
s <= std_logic_vector(unsigned(a) + unsigned(b));
\end{code}
\end{bloc}

\begin{bloc}{Machine à états (Moore)}{ex}
\begin{code}
type etat_t is (IDLE, ATT, RUN);
signal etat : etat_t;
process (clk)
begin
  if rising_edge(clk) then
    case etat is
      when IDLE => if go = '1' then etat <= ATT; end if;
      when ATT => if ok = '1' then etat <= RUN; end if;
      when RUN => if fin = '1' then etat <= IDLE; end if;
    end case;
    if rst = '1' then  -- reset synchrone : la
      etat <= IDLE;    -- dernière affectation
    end if;            -- l'emporte
  end if;
end process;
busy <= '1' when etat = RUN else '0';  -- sorties
\end{code}
\end{bloc}

\begin{bloc}{Trois états}{ex}
\begin{code}
b   <= a when oe = '1' else (others => 'Z');  -- sortie
din <= pad;                                   -- inout lu
\end{code}
\end{bloc}

\begin{bloc}{Afficher des valeurs}{io}
\begin{code}
report "cnt = " & integer'image(to_integer(cnt));
report "v = " & to_hstring(v);  -- VHDL-2008 (hexa)
report "t = " & time'image(now);
\end{code}
\end{bloc}
\end{colonne}\colbreak
% ======================= colonne 3 : simulation =======================
\begin{colonne}
\begin{bloc}{Testbench}{io}
\begin{code}
library ieee;
use ieee.std_logic_1164.all;
entity tb_and2 is             -- pas de ports
end entity;
architecture sim of tb_and2 is
  signal a, b, c : std_logic := '0';
  signal clk     : std_logic := '0';
begin
  clk <= not clk after 5 ns;  -- horloge 100 MHz
  uut : entity work.and2
    port map (a => a, b => b, c => c);
  stim : process
  begin
    a <= '1'; b <= '1';
    wait for 10 ns;
    assert c = '1' report "1 and 1 doit donner 1"
      severity error;
    wait until rising_edge(clk);
    report "fin" severity note;
    std.env.stop;             -- VHDL-2008 ; sinon wait;
  end process;
end architecture;
\end{code}
\begin{optab}[20mm]
\V|wait for t| & attendre une durée \\
\V|wait until c| & attendre que \V|c| devienne vrai \\
\V|wait on s| & attendre un changement de \V|s| \\
\V|wait;| & suspendre pour toujours \\
\V|severity| & \V|note warning error failure| (arrêt)
\end{optab}
\end{bloc}

\begin{bloc}{RAM synchrone}{ex}
\begin{code}
type ram_t is array (0 to 2**AW - 1)
  of std_logic_vector(DW-1 downto 0);
signal ram : ram_t;
process (clk)
begin
  if rising_edge(clk) then
    if we = '1' then
      ram(to_integer(unsigned(adr))) <= din;
    end if;
    dout <= ram(to_integer(unsigned(adr)));  -- 1 cycle
  end if;
end process;
\end{code}
\note{Motif reconnu comme bloc mémoire (BRAM) à la synthèse.}
\end{bloc}

\begin{bloc}{Pièges}{ex}
\begin{optab}[19mm]
Signal / variable & lire un signal juste après l'avoir affecté rend l'ancienne valeur. \\
Latch & un \V|if| sans \V|else| dans un process combinatoire mémorise : mettre une valeur par défaut. \\
Pilotes multiples & deux process sur un signal : \V|'X'| si les valeurs s'opposent, refusé sur \V|integer|. \\
Mots réservés & \V|bus|, \V|in|, \V|out|, \V|range|... ne peuvent pas nommer un signal. \\
\V|downto| / \V|to| & \V|(7 downto 0)| : bit 7 à gauche (MSB) ; ne pas mélanger. \\
Init. & \V|signal s : ... := '0'| ne sert qu'en simulation (ou FPGA) : prévoir un reset.
\end{optab}
\end{bloc}
\end{colonne}
\end{document}
